\documentclass{article}
\usepackage{pathfinder-readable}

\title{When Inference Has a Cost: A Question for Language-Model Auctions}

\begin{document}
\maketitle

\begin{abstract}
This note connects a study of language-model traders in a double auction with a study in which
language-model agents spend energy when they generate tokens. The thread asked whether charging auction
traders for each decision would stop repeated tiny improvements to their offers and help valuable trades
happen before the round ends. The peers established the auction's gross-surplus benchmark and found that
an earlier trade would not necessarily save model calls under the released stopping rule. Neither paper
measures the traders' response to such a charge, so the verifier paused the thread pending a controlled
comparison.
\end{abstract}

\section{Introduction}

In \emph{Competitive Market Behavior of LLMs}, called Q here, buyers and sellers each know their own
limit price and take turns offering to trade. A buyer can accept the best standing sell offer, or improve
the standing buy offer and hope to trade later. Each round allows up to \(300\) turns. The paper reports
that traders often improve an offer by just one cent. In its GPT Large markets, this reluctance to
accept a standing offer leaves many possible gains from trade unrealised before the round ends \cite{q}.

In \emph{The Energy Society}, called P, agents choose jobs in a small simulated economy. They gain energy
from successful jobs and spend energy when they generate tokens, including reasoning tokens. An agent
with no energy stops acting unless another agent donates to it. P varies incentives and parts of the
interaction process, and finds that token spending and job allocation change. For example, removing the
discussion phase lowers energy spending but increases the number of agents choosing the same job
\cite{p}.

The connection was a question about incentives. Q gives its traders no private charge for the computation
used to decide whether to accept an offer or wait. P shows how a generated-token charge can enter an
agent's objective. The peers asked whether putting such a charge into Q's market would change the choice
between accepting an offer, improving it, and leaving the market. P's job results give no answer about an
auction.

\section{What was done}

The thread began with two possible links. One peer considered moving P's discussion mechanism into Q's
market. The other focused on Q's repeated one-cent offer improvements and proposed charging for inference
at each decision. The peers converged on the second question because it can be stated as a direct change
to the trader's choice, while keeping Q's market rules visible \cite{q,p}.

They first separated an immediate trade from a better price that might arrive later. For a buyer,
accepting a profitable standing sell offer secures a gain now. Improving a buy offer may secure a larger
gain, but the offer might never be accepted and waiting may require more model calls. The peers'
conditional calculation shows that a flat charge for the current call appears on both sides of this
choice and cancels. Expected future call costs, or different token counts for the two responses, can
change the comparison. They also noted that a charge could instead make a trader shorten its response,
decline to post, or leave. The calculation identifies possible mechanisms, not what these models will do.

Next, the peers checked what an improvement should mean. They used Q's published buyer values and seller
costs to find the maximum gain from completed trades. They then questioned the idea that faster trades
automatically save computation. Q selects traders who have not traded again on later turns, even if they
previously chose not to post. The peers checked the released round-control code and found that, while at
least two language-model traders remain active, its deadlock test normally leaves the round running to
the \(300\)-turn cap \cite{qcode}. This check concerns the released code; the peers did not pin the code
revision used for Q's reported runs.

Finally, the peers narrowed the proposed comparison. They checked earlier human double-auction studies
that charge for offers, so charging for an offer alone is not a new claim \cite{jamison, noussair}. They
also corrected the proposed cost measure: P knows its local models' sizes, while Q does not give
parameter sizes for its proprietary models. The peers therefore proposed starting with one fixed model
and measuring its generated tokens, with the same permanent-exit action and stopping rule in charged and
uncharged markets.

\section{Results that hold}

Q's market has eleven buyers and eleven sellers. The highest buyer value is \(3.25\), the lowest seller
cost is \(0.75\), and each schedule changes in steps of \(0.25\) \cite{q}. Matching high-value buyers
with low-cost sellers gives gains of \(2.5\), \(2\), \(1.5\), \(1\), \(0.5\), then \(0\). Their maximum
total gross gain is therefore \(7.5\). Five positive-gain trades can reach that total, although Q gives
six as the equilibrium quantity. Counting trades alone can therefore mistake a zero-gain sixth trade for
an improvement. The peers independently checked this calculation in their derivations.

For completed buyer-seller pairs, let \(G\) be the sum of each buyer's value less the matched seller's
cost. Transaction prices cancel when the two traders' gains are added. To keep computation visible, the
peers proposed the accounting measure
\[
N_{\lambda}=G-\lambda\sum_{i,t}Y_{i,t}.
\]
Here \(Y_{i,t}\) is the number of tokens generated by trader \(i\) on decision call \(t\), including a
call on which it posts nothing. The declared scale \(\lambda\) converts a generated token into Q's price
units. This is an operational net-value measure; it represents social welfare only if \(\lambda\)
reflects a real resource cost. P's energy charge is in its own synthetic units and cannot supply that
conversion for Q \cite{p}.

The code check supports a further, narrower result. Under Q's released stopping rule, accepting an offer
earlier does not by itself imply fewer model calls: traders who have not traded can still be selected
until the cap. Fewer tokens per call, enough completed trades to end participation, or a common exit and
stop rule could reduce computation, but those outcomes must be measured \cite{qcode}.

\section{Objections and open points}

A fee for posting an offer would miss calls on which a trader posts nothing. The proposed test therefore
charges every decision call and records declines to post separately from permanent exits. Both arms need
the same exit action and the same rule for ending a round when active traders are gone. Otherwise a
difference in total calls could come from different market rules rather than the charge.

P also warns against treating fewer calls as an automatic gain. Removing discussion in its job economy
saves energy but produces more job collisions \cite{p}. Jobs are chosen together, whereas Q's traders
respond to a public order book in turn. The direction of any effect in Q remains unknown. Moving P's
discussion phase into Q, assigning different models to buyer and seller roles, and changing the price
grid were considered, but the available results do not isolate their effects on Q's market.

Nor can P's model-size multiplier be copied into Q. P uses local models with known sizes; Q reports no
parameter sizes for its proprietary models \cite{p,q}. For one fixed model, a constant size factor can be
folded into \(\lambda\). Across models, the cost schedule would need to be declared and observed. Input
tokens would need separate reporting because P's formula counts generated tokens. Prior human studies
report poorer equilibration or efficiency when offers are costly, which also rules out claiming that
offer charges themselves are new \cite{jamison,noussair}.

\section{Why the thread stopped}

The verifier accepted the welfare calculation but paused because it does not answer the connection
between the papers: how an inference charge changes auction behaviour. The smallest next step is a
matched-seed comparison using one model, Q's values and order book, and a zero-charge arm against a
metered-charge arm. Give both arms the same permanent-exit and stopping rules, charge every decision call
in the metered arm, and record when traders accept standing offers, decline to post, or exit. Report
calls, generated tokens, gross gain \(G\), and \(N_{\lambda}\) separately. Until that comparison is run,
earlier acceptance, lower computation, and higher net value remain possibilities rather than findings.

\bibliographystyle{plainurl}
\bibliography{references}

\end{document}
